\documentclass{article}
\usepackage[utf8]{inputenc}
\usepackage{booktabs}
\usepackage{siunitx}
\usepackage{caption}
\usepackage{threeparttable}

\sisetup{
    separate-uncertainty=true,
    table-align-uncertainty=true,
    table-format=1.3,
    tight-spacing=true,
    input-symbols = { ( ) },
    table-align-text-after=false
}

\begin{document}

\begin{table}[htbp]
    \centering
    \begin{threeparttable}
    \caption{The Multi-faceted Impact of Policy $T$: Main Effects vs. Mechanisms}
    \label{tab:panel_story}
    \begin{tabular}{l *{3}{S}}
        \toprule
        & {(1)} & {(2)} & {(3)} \\
        \midrule
        \multicolumn{4}{l}{\textbf{Panel A: Main Treatment Effects}} \\
        Treatment ($D$) & 0.850*** & 0.825*** & 0.790*** \\
                        & (0.045) & (0.042) & (0.041) \\
        \addlinespace
        R-squared       & 0.245 & 0.289 & 0.312 \\
        \midrule
        \multicolumn{4}{l}{\textbf{Panel B: Possible Mechanisms (Intermediate Outcomes)}} \\
        & {Labor Supply} & {Human Capital} & {Efficiency} \\
        \cmidrule(lr){2-4}
        Treatment ($D$) & 0.450*** & 0.125** & 0.089 \\
                        & (0.120) & (0.055) & (0.065) \\
        \addlinespace
        R-squared       & 0.185 & 0.145 & 0.056 \\
        \midrule
        Controls        & {Yes} & {Yes} & {Yes} \\
        Fixed Effects   & {Year} & {State} & {Both} \\
        Observations    & {45,200} & {45,200} & {45,200} \\
        \bottomrule
    \end{tabular}
    \begin{tablenotes}
        \small
        \item \textit{Notes:} Panel A reports the intent-to-treat effect on the primary outcome. Panel B explores potential transmission channels by using intermediate variables as dependent variables. Total sample size is consistent across panels.
        \item *** p<0.01, ** p<0.05, * p<0.1.
    \end{tablenotes}
    \end{threeparttable}
\end{table}

\end{document}
